\section{Introduction}

Hotel reviews on platforms such as TripAdvisor are widely used by travellers and
hotel operators alike. Reviews with four or five stars are generally assumed to
express satisfaction, yet they frequently contain concrete suggestions for
improvement embedded within otherwise positive text. A guest who rates a hotel
five stars may still write, ``I wish the breakfast had more vegetarian
options.'' Standard sentiment analysis classifies such reviews as positive,
missing the constructive content within them.

Suggestion mining is the task of identifying sentences that propose an
improvement, give advice, or recommend an action \citep{negi2018open}. Framed as
sentence-level binary classification, it distinguishes suggestions from
non-suggestions. The task was formalised in the SemEval-2019 Task~9 shared task
\citep{negi2019semeval}, which provided labelled data from software developer
forums. That shared task demonstrated that BERT-based models achieved the best
performance for in-domain suggestion detection but also revealed that
cross-domain transfer---applying a model trained on one domain to
another---remains a challenge \citep{park2019bert}.

While SemEval-2019 Task~9 Subtask~B evaluated cross-domain
transfer to hotel reviews, no prior work combines
domain-adapted annotation guidelines, a multi-tier model
comparison, and thematic analysis for suggestion mining in
high-rated hotel reviews. Hotel reviews differ from software forums in
vocabulary, sentence structure, and the way suggestions are expressed:
software forum users write explicit feature requests (``The API should support
batch uploads''), while hotel guests embed suggestions within narrative
descriptions (``The room was lovely but the pillows could be firmer''). This
domain shift makes direct transfer from existing datasets non-trivial, and
raises the question of how much domain-specific adaptation is needed.

Our research question is: \textbf{To what extent can NLP models detect
actionable improvement suggestions within high-star hotel reviews?}

We address this question through four contributions:
\begin{itemize}
  \item \textbf{Domain-adapted annotation guidelines} for identifying
  suggestions in hotel reviews, with inter-annotator agreement measured on a
  calibration sample of 100 sentences.
  \item \textbf{A three-tier model comparison}---regex baseline, TF-IDF with
  logistic regression (TF-IDF~+~LR) baseline, and BERT classifier---that
  quantifies the value of increasingly rich representations for suggestion
  detection.
  \item \textbf{Two-stage fine-tuning} of BERT: first on the SemEval-2019
  training data, then on annotated hotel review sentences, with a cross-domain
  evaluation that measures the performance delta between in-domain and
  cross-domain settings.
  \item \textbf{BERTopic-based thematic analysis} of extracted suggestions,
  grouping them into hotel-relevant categories (room, food and beverage, service,
  facilities, location, value).
\end{itemize}
